\section{The Shape of the Model: A Random Hyperbolic Network}
\label{sec:shape}

Before the formal definitions, the picture in plain terms, since the whole framework
rests on one unfamiliar kind of object, and the natural guesses about it are all wrong
in instructive ways.

\paragraph{What it is.}
The mathematical structure is a graph that grows. Picture a vast web of points, each
point a vibe, joined by links, each link a note. The web is not a grid. Its points are
scattered at random, and the space they are scattered in is hyperbolic, negatively
curved, the saddle-shaped geometry in which the room around any point keeps opening up
the further you look, so the number of points within reach grows exponentially with
distance. On this web sits a simple rule that resets the state of each point from its
neighbours, run over and over, while the web itself keeps adding points at its leading
edge. In the language of discrete physics this is a causal set carrying a network
automaton: a random, hyperbolic, growing graph with a local update rule. That object,
and nothing richer, is the model.

\paragraph{How it works.}
Each point carries a single three-valued state, its tone, one of $\{-1,0,+1\}$, read as
pain, peace, and pleasure. It is joined to about ten others, not a fixed number but a
spread, the ones it directly experiences. Each link carries its own three-valued number,
the fill, the strength and sign of that attention. To update, a point reads the outward
states (the wills) of its neighbours, multiplies each by the fill of the link, adds the
results into a single whole number, and takes the sign of that sum as its next state: a
positive total becomes pleasure, a negative one pain, a tie peace. There are no
fractional weights and no stored real numbers anywhere, only three-valued states and a
transient integer count, exactly the spirit of a cellular automaton's tally of its
neighbours. The points update one at a time with no global clock, and the web only ever
grows, its links accumulating and never undone.

\paragraph{What it is not.}
It is not a regular cellular automaton on a square or cubic grid. A grid has preferred
directions, so on a grid the speed of light would depend on direction and Lorentz
invariance would break, which experiment forbids. It is not a regular hyperbolic tiling
either, the orderly Escher-like packings of the hyperbolic plane, for the same reason,
they are still too rigid. The randomness is not a convenience but the load-bearing
choice, since only a random scatter has no preferred frame even at the smallest scale.
It is not a field theory on a smooth space, because there is no smooth space underneath.
And there are no hidden variables behind the tones, and no continuous amplitudes stored
on any point.

\paragraph{Where the smooth physics lives.}
If nothing continuous is stored, where do the smooth numbers of physics come from? They
are large-scale summaries of many points, the way a temperature summarises jiggling
molecules. The geometry of spacetime is the shape of the web seen from far off. Energy
and the flow of time are what a local operator, the graph's Laplacian, does on the grown
web. A quantum amplitude is a patch of many points read together, its magnitude the
fraction of them excited, its phase their shared rhythm. Each of these is real and
measurable, and none is written on any single vibe. The base is discrete and
three-valued, the continuum is the view from a distance.

This is the object the next section defines precisely and the rest of the paper tests: a
growing, random, hyperbolic network of three-valued points that vote locally, with all
the smooth structure of the world appearing only in the large.
